\section{Discussion}
\label{sec:limits}

The decomposition separates two benchmarks that a headline accuracy figure does not. On HC3, orthography alone matches content alone, 3.20\% against 3.26\% test error, and a paired test over 10,732 documents cannot distinguish that 0.07-point difference from zero, nor can the same test on any of five independent partitions. A fine-tuned DeBERTa improves on either arm by roughly 0.03 weighted F1. On DAIGT V2, content is stronger than orthography by a factor of 7.9 in error. The best classical model there is within 0.0007 weighted F1 of the best transformer, a gap the same test also cannot separate from zero, though only because the two models read different amounts of text. A practitioner choosing an evaluation corpus, or a reviewer reading a reported score, cannot recover that distinction from the score itself.

The attribution in Section~\ref{sec:shap} says where the difference comes from. HC3's surface arm is close to a single-feature model, since the rate of spaces before punctuation leads the next feature by a factor of 2.3 and the corpus carries that cue in 88.7\% of human documents against 0.3\% of machine ones. DAIGT V2's surface arm reads a diffuse mixture of length, word shape and casing with no comparable leader. Those are different kinds of corpus even though both are separable by form, and the practical consequence differs accordingly. A single collection convention can be cleaned. A diffuse stylistic signal cannot, and it may not even be an artefact.

What the measurement does not license is a claim that either corpus is clean. DAIGT V2's surface-only arm reaches 0.9214 weighted F1, far above chance, so surface form is substantially informative on both. The finding is comparative, that surface reaches parity with content on one corpus and not the other, and it is bounded by the arms we built. A richer surface featurisation would raise both surface numbers and a stronger content model would raise both content numbers, so the arms bound rather than partition the separability, in the sense made explicit in Section~\ref{sec:decomp-def}.

The two arms are also not disjoint as built, since document length reaches both. Closing that channel weakens neither conclusion. DAIGT V2's content advantage widens from 7.9 to 10.6 times, and HC3's parity becomes a 2.91-point advantage for orthography, because removing length costs the content arm eighteen times what it costs the surface arm. The control nearly produced a third false conclusion of its own, since normalising bag-of-words rows without restoring their scale collapses the arm through regularisation and reads as a length effect.

Two results here began as the opposite conclusion, and both are worth stating for what they say about method rather than about these corpora. A cleaning experiment that reports no change is uninterpretable without evidence that the pipeline ran at all, and a tokeniser that cannot represent the cue being cleaned makes the null a fact about the model rather than about the corpus. A collapse under perturbation is uninterpretable without a label-free control, since a model that calls random strings human at 0.98 confidence will produce the same curve with no vulnerability behind it. We report both requirements rather than the findings they dissolve, and we suggest comparable published collapses be read with the second control in mind.

There is a connection to fairness that we did not measure and that seems to us the most consequential reason to care which arm a score comes from. Liang et al.~\cite{Liang2023gpt} show deployed detectors misclassify non-native English writing as machine-generated at high rates. A detector whose separability rests substantially on orthographic regularity is a detector whose false positives should be expected to concentrate on writers whose orthographic habits differ from the corpus it was fitted to. Our measurement predicts which corpora carry that risk most sharply. It does not demonstrate the effect, and establishing it would need annotated author metadata that neither corpus provides.

\section{Limitations}

Several limits bound every number above, and we state them rather than leaving them to be inferred.

The transformers see at most 128 tokens, so the DAIGT V2 transformer results describe classification from an essay's opening rather than from the essay. We measure that rather than assume it, by refitting the classical grid on the same window, and it reverses the classical-versus-transformer reading on that corpus. The cross-corpus comparison in Table~\ref{tab:cross} remains not like-for-like in the amount of text each model receives.

Seed spread is estimated from three seeds, and one such estimate moved by 31\% when its runs were repeated, because GPU training here is not bitwise deterministic and best-epoch selection can choose a materially different checkpoint on a rerun. No claim in this paper rests on a seed range, and every comparison uses a paired test on saved predictions instead.

The content arm keeps stopwords and inflections, so it is not a reduced content model. It is still a bag-of-words one, and word order and syntax are outside it, so it bounds content-carried separability from below for that reason rather than through its preprocessing. A syntactic or contextual content arm would raise $\varepsilon_{\mathrm{C}}$'s competitor and could narrow the DAIGT V2 gap.

The surface set is 47 hand-built features, two of which have no variance on the DAIGT V2 test split and one of which has none on HC3. It is deliberately cheap and it is certainly not exhaustive, so $\varepsilon_{\mathrm{S}}$ is an upper bound on surface-arm error rather than a measurement of everything surface form could carry.

HC3 has one generator collected at one time, and both corpora are English. We evaluate two corpora and five model families, and nothing here establishes that the pattern holds beyond them. In particular we do not claim that HC3 is representative of short-form question answering, or DAIGT V2 of student writing, only that these two widely used benchmarks differ from one another on the axis we measure.

\section{Conclusion}

We proposed a decomposition that measures how much of a machine-generated-text benchmark's separability is available from surface form alone rather than from content, and applied it to DAIGT V2 and HC3 across five model families and eight configurations per corpus. On HC3 the two arms are indistinguishable under paired testing on five independent partitions, and a per-feature attribution shows a single collection convention accounts for most of the surface arm. On DAIGT V2 content leads by a factor of 7.9 in error, and the surface signal is diffuse rather than concentrated. Closing the document-length channel that both arms share strengthens both conclusions. Three controls, on tokenisation, on pipeline execution and on label-free inputs, each overturned a conclusion that this study had reached without them, and we report the controls as requirements rather than the findings they dissolved.

The measurement itself is the portable part. It needs no new annotation, it fits two logistic regressions, and it runs in minutes on a laptop. It yields one number per corpus that a headline accuracy cannot express, namely how much of that corpus a model could pass without reading the language. We would encourage reporting it alongside any new detection benchmark, in the same way a hypothesis-only baseline is now reported alongside a natural language inference dataset~\cite{Poliak2018hypothesis}.
